\textbf{Phần 1. Trắc nghiệm lựa chọn}
\setcounterpageref{ex}{0}

\begin{ex} %[3D2N10]
Kết quả của phép nhân $7 \times 3$ là
\choice
{$18$}
{$24$}
{$28$}
{\True $21$}
\loigiai{
$7 \times 3 = 21$.\\
\textbf{Đáp án của câu hỏi là D}
}
\end{ex}

\begin{ex} %[3D2N10]
Kết quả của phép chia $28 : 7$ là
\choice
{\True $4$}
{$3$}
{$5$}
{$6$}
\loigiai{
$28 : 7 = 4$.\\
\textbf{Đáp án của câu hỏi là A}
}
\end{ex}

\begin{ex} %[3D2H10]
Kết quả của phép nhân $7 \times 8$ là
\choice
{$54$}
{$58$}
{$63$}
{\True $56$}
\loigiai{
$7 \times 8 = 56$.\\
\textbf{Đáp án của câu hỏi là D}
}
\end{ex}

\begin{ex} %[3D2H10]
Kết quả của phép chia $63 : 7$ là
\choice
{$8$}
{\True $9$}
{$7$}
{$6$}
\loigiai{
$63 : 7 = 9$.\\
\textbf{Đáp án của câu hỏi là B}
}
\end{ex}

\begin{ex} %[3D2H10]
Điền số thích hợp vào chỗ chấm: $7 \times \dots = 49$.
\choice
{$6$}
{$8$}
{\True $7$}
{$9$}
\loigiai{
$49 : 7 = 7$.\\
\textbf{Đáp án của câu hỏi là C}
}
\end{ex}

\begin{ex} %[3D2H10]
Điền số thích hợp vào chỗ chấm: $\dots : 7 = 5$.
\choice
{$30$}
{\True $35$}
{$40$}
{$42$}
\loigiai{
$5 \times 7 = 35$.\\
\textbf{Đáp án của câu hỏi là B}
}
\end{ex}

\begin{ex} %[3D2V10]
Tính giá trị biểu thức: $7 \times 6 + 18$.
\choice
{$58$}
{$62$}
{\True $60$}
{$64$}
\loigiai{
$7 \times 6 + 18 = 42 + 18 = 60$.\\
\textbf{Đáp án của câu hỏi là C}
}
\end{ex}

\begin{ex} %[3D2V10]
Một tuần lễ có 7 ngày. Hỏi 5 tuần lễ có bao nhiêu ngày?
\choice
{\True $35$ ngày}
{$30$ ngày}
{$40$ ngày}
{$28$ ngày}
\loigiai{
$7 \times 5 = 35$ ngày.\\
\textbf{Đáp án của câu hỏi là A}
}
\end{ex}

\begin{ex} %[3D2V10]
Có 42 học sinh xếp đều thành 7 hàng. Mỗi hàng có số học sinh là
\choice
{$5$ học sinh}
{$7$ học sinh}
{$8$ học sinh}
{\True $6$ học sinh}
\loigiai{
$42 : 7 = 6$ học sinh.\\
\textbf{Đáp án của câu hỏi là D}
}
\end{ex}

\begin{ex} %[3D2V10]
Mỗi gói bánh giá $7,000$ đồng. Mua 4 gói bánh như thế hết số tiền là
\choice
{$24,000$ đồng}
{$32,000$ đồng}
{\True $28,000$ đồng}
{$21,000$ đồng}
\loigiai{
$7,000 \times 4 = 28,000$ đồng.\\
\textbf{Đáp án của câu hỏi là C}
}
\end{ex}

\begin{ex} %[3D2V10]
Có 56 quyển vở thưởng đều cho 7 bạn học sinh giỏi. Mỗi bạn nhận được số quyển vở là
\choice
{\True $8$ quyển}
{$7$ quyển}
{$9$ quyển}
{$6$ quyển}
\loigiai{
$56 : 7 = 8$ quyển vở.\\
\textbf{Đáp án của câu hỏi là A}
}
\end{ex}

\begin{ex} %[3D2C10]
Tìm một số biết rằng nếu lấy số đó chia cho 7 rồi cộng với 15 thì được kết quả là 21.
\choice
{$35$}
{\True $42$}
{$49$}
{$28$}
\loigiai{
Số đó chia cho 7 bằng: $21 - 15 = 6$.\\
Số đó là: $6 \times 7 = 42$.\\
\textbf{Đáp án của câu hỏi là B}
}
\end{ex}

\textbf{Phần 2. Trắc nghiệm đúng sai}
\setcounterpageref{ex}{0}

\begin{ex} %[3D2H10]
Xét tính đúng sai của các phép tính trong bảng nhân 7, bảng chia 7:
\choiceTF
{\True $7 \times 4 = 28$.}
{\True $35 : 7 = 5$.}
{$7 \times 9 = 65$.}
{\True $7 \times 9 = 63$.}
\loigiai{
\begin{itemize}
    \item $7 \times 4 = 28$: Đúng.
    \item $35 : 7 = 5$: Đúng.
    \item $7 \times 9 = 63 \neq 65$. Ý 3 Sai, Ý 4 Đúng.
\end{itemize}
\begin{itemize}
    \item a) \textbf{Đúng}.
    \item b) \textbf{Đúng}.
    \item c) \textbf{Sai}.
    \item d) \textbf{Đúng}.
\end{itemize}
}
\end{ex}

\begin{ex} %[3D2V10]
Có 7 hàng cây, mỗi hàng có 7 cây. Người ta trồng thêm 11 cây nữa. Xét tính đúng sai:
\choiceTF
{\True Lúc đầu 7 hàng có tất cả 49 cây.}
{\True Sau khi trồng thêm, tổng số cây là 60 cây.}
{Số cây trồng thêm nhiều hơn số cây trong 1 hàng là 5 cây.}
{\True Số cây trồng thêm nhiều hơn số cây trong 1 hàng là 4 cây.}
\loigiai{
\begin{itemize}
    \item Lúc đầu: $7 \times 7 = 49$ cây. Đúng.
    \item Tổng: $49 + 11 = 60$ cây. Đúng.
    \item Hiệu: $11 - 7 = 4$ cây $\neq 5$. Ý 3 Sai, Ý 4 Đúng.
\end{itemize}
\begin{itemize}
    \item a) \textbf{Đúng}.
    \item b) \textbf{Đúng}.
    \item c) \textbf{Sai}.
    \item d) \textbf{Đúng}.
\end{itemize}
}
\end{ex}

\textbf{Phần 3. Trả lời ngắn}
\setcounterpageref{ex}{0}

\begin{ex} %[3D2H10]
Tính tích: $7 \times 7$.
\shortans{49}
\loigiai{
\textbf{Đáp số:} $49$\\
$7 \times 7 = 49$.
}
\end{ex}

\begin{ex} %[3D2H10]
Tính thương: $21 : 7$.
\shortans{3}
\loigiai{
\textbf{Đáp số:} $3$\\
$21 : 7 = 3$.
}
\end{ex}

\begin{ex} %[3D2V10]
Tính giá trị biểu thức: $7 \times 8 - 26$.
\shortans{30}
\loigiai{
\textbf{Đáp số:} $30$\\
$56 - 26 = 30$.
}
\end{ex}

\begin{ex} %[3D2V10]
Một đàn gà có 7 chuồng, mỗi chuồng có 8 con gà. Hỏi có tất cả bao nhiêu con gà?
\shortans{56}
\loigiai{
\textbf{Đáp số:} $56$\\
$8 \times 7 = 56$ con gà.
}
\end{ex}

\begin{ex} %[3D2V10]
Có 70 lít dầu chia đều vào 7 thùng. Hỏi mỗi thùng có bao nhiêu lít dầu?
\shortans{10}
\loigiai{
\textbf{Đáp số:} $10$\\
$70 : 7 = 10\text{ l}$.
}
\end{ex}

\begin{ex} %[3D2C10]
Tìm $x$ biết: $x \times 7 = 40 + 23$.
\shortans{9}
\loigiai{
\textbf{Đáp số:} $9$\\
$x \times 7 = 63 \Rightarrow x = 63 : 7 = 9$.
}
\end{ex}
